\documentclass[10pt]{article}
\usepackage[T1]{fontenc}
\usepackage[utf8]{inputenc}
\usepackage[french]{babel}
\usepackage{amsmath,amssymb,mathrsfs,mathpazo}
\usepackage{geometry}
\geometry{a4paper,margin=2cm}
\usepackage{graphicx,array,tabularx,tabularray,colortbl,multirow,multicol,booktabs,enumitem,pifont,ulem,fancybox,diagbox,cancel}
\usepackage{pstricks,pstricks-add,pst-plot,pst-tree,pst-node,pst-text,pst-func,pst-math,pst-bspline,pst-3dplot,pst-eucl,pst-solides3d}
\let\pscalculate\undefined
\usepackage{pst-bezier}
\usepackage{tikz,pgfplots,tkz-tab}
\usetikzlibrary{arrows,arrows.meta,patterns,calc,positioning,trees,decorations.markings,shapes.geometric,shapes.misc,shapes.arrows,intersections,decorations.pathmorphing,decorations.pathreplacing,backgrounds,shadows,babel}
\pgfplotsset{compat=1.18}
\usepgfplotslibrary{fillbetween}
\usepackage[np]{numprint}
\usepackage{listings,makecell,pas-tableur,twoopt}
\let\barre\undefined
\usepackage[pstricks]{tablvar}
\definecolor{darkgreen}{RGB}{0,100,0}
\providecommand{\red}{\color{red}}
\providecommand{\blue}{\color{blue}}
\providecommand{\green}{\color{green}}
\providecommand{\black}{\color{black}}
\providecommand{\white}{\color{white}}
\providecommand{\gray}{\color{gray}}
\providecommand{\cyan}{\color{cyan}}
\providecommand{\magenta}{\color{magenta}}
\providecommand{\yellow}{\color{yellow}}
\providecommand{\violet}{\color{violet}}
\providecommand{\orange}{\color{orange}}
\providecommand{\euro}{\texteuro}
\providecommand{\EUR}{\texteuro}
\providecommand{\ieme}{\textsuperscript{e}}
\providecommand{\degres}{\ensuremath{^\circ}}
\providecommand{\degree}{\ensuremath{^\circ}}
\providecommand{\vv}[1]{\overrightarrow{#1}}
\providecommand{\Coord}[1]{#1}
\providecommand{\decosix}{}
\providecommand{\starredbullet}{\ensuremath{\bullet}}
\pagestyle{empty}
\setlength{\parindent}{0pt}
\usepackage{adjustbox}\DeclareUnicodeCharacter{2236}{:}\DeclareUnicodeCharacter{2190}{\ensuremath{\leftarrow}}
\providecommand{\R}{\mathbb{R}}
\providecommand{\N}{\mathbb{N}}
\providecommand{\Z}{\mathbb{Z}}
\providecommand{\Q}{\mathbb{Q}}
\providecommand{\D}{\mathbb{D}}
\providecommand{\C}{\mathbb{C}}
\providecommand{\e}{{\mathrm{e}}}
\providecommand{\vect}[1]{\overrightarrow{#1}}
\providecommand{\vectt}[1]{\overrightarrow{\mathrm{#1}}}
\providecommand{\vvt}[1]{\overrightarrow{\mathrm{#1}}}
\providecommand{\np}[2][]{{#2~#1}}
\providecommand{\vv}[1]{{\overrightarrow{#1}}}
\renewcommand{\d}{{\mathrm{d}}}
\providecommand{\degres}{{^\circ}}
\providecommand{\degre}{{^\circ}}
\providecommand{\blue}{\color{blue}}
\providecommand{\red}{\color{red}}
\providecommand{\enskip}{\quad}
\providecommand{\llbracket}{\lbrack\!\lbrack}
\providecommand{\rrbracket}{\rbrack\!\rbrack}
\providecommand{\euro}{{\text{€}}}
\providecommand{\no}{{\text{n°}}}
\providecommand{\interval}[2]{{[#1;#2]}}
\providecommand{\intervFO}[2]{{[#1;#2[}}
\providecommand{\dotfill}{\ldots}

\providecommand{\ds}{\displaystyle}
\providecommand{\textding}[1]{#1}
\providecommand{\Oij}{\(\mathrm{O};\vec{\imath},\vec{\jmath}\)}
\providecommand{\Oijk}{\(\mathrm{O};\vec{\imath},\vec{\jmath},\vec{k}\)}
\providecommand{\Ouv}{\(\mathrm{O};\vec{u},\vec{v}\)}
\providecommand{\intervFF}[2]{\ensuremath{[#1;#2]}}
\renewcommand{\euro}{\ifmmode\text{\texteuro}\else\texteuro\fi}
\DeclareUnicodeCharacter{00B0}{\ifmmode^{\circ}\else\textdegree\fi}
\DeclareUnicodeCharacter{20AC}{\ifmmode\text{\texteuro}\else\texteuro\fi}
% Énoncés : Éducation nationale. Transcription et corrigés : APMEP, Denis Vergès et contributeurs. Découpage et classement : CoopMaths, Sébastien Lozano, Éric Elter et contributeurs. Adaptation de la présentation : Maths Explorer.
% CC BY-SA : https://coopmaths.fr/www/about/
% https://coopmaths.fr/alea/static/bac/2023/tex/bac_2023_03_sujet1_asie_3.tex
\begin{document}
\noindent\textit{APMEP / CoopMaths. CC BY-SA. Présentation adaptée par Maths Explorer.}\par\medskip
\medskip

Soit $k$ un réel strictement positif.

Le but de cet exercice est de déterminer le nombre de solutions de l'équation 
\[\ln (x) = kx\]
 de paramètre $k$.
 
\smallskip

\begin{enumerate}
\item \textbf{Conjectures graphiques :}

\medskip

On a représenté, ci-dessous, dans un repère orthogonal, la courbe d'équation $y = \ln (x)$, la droite d'équation $y = x$ ainsi que la droite d'équation $y = 0,2x$ :

\begin{center}
\psset{unit=0.8cm,arrowsize=2pt 3}
\begin{pspicture*}(-1,-1.2)(16.5,9)
\psaxes[linewidth=1.25pt]{->}(0,0)(0,-1)(16.5,9)
\psplot[plotpoints=2000,linewidth=1.25pt,linecolor=red]{0.1}{16.5}{x ln}\uput[d](15,2.3){\red $y = \ln x$}
\psplot[plotpoints=2000,linewidth=1.25pt,linecolor=blue]{0}{16.5}{x 0.2 mul}\uput[4](14,3.3){\blue $y = 0,2 x$}
\psplot[plotpoints=2000,linewidth=1.25pt]{0}{16.5}{x}\uput[ul](8.5,8.5){$y = x$}
\end{pspicture*}
\end{center}

À partir du graphique ci-dessus, conjecturer le nombre de solutions de l'équation $\ln (x) = kx$ pour $k = 1$ puis pour $k = 0,2$.
\item Étude du cas $k = 1$ :

On considère la fonction $f$, définie et dérivable sur $]0~;~+\infty[$, par: 
\[f(x) = \ln (x) - x.\]

On note $f'$ la fonction dérivée de la fonction $f$.
	\begin{enumerate}
		\item Calculer $f'(x)$.
		\item Étudier le sens de variation de la fonction $f$ sur $]0~;~+\infty[$.

Dresser le tableau des variations de la fonction $f$ en y faisant figurer la valeur exacte des extremums s'il y en a.

Les limites aux bornes de l'intervalle de définition ne sont pas attendues.
		\item En déduire le nombre de solutions de l'équation $\ln (x) = x$.
	\end{enumerate}	
\item Étude du cas général :

$k$ est un nombre réel strictement positif.

On considère la fonction $g$ définie sur $]0~;~+\infty[$ par:
\[g(x) = \ln (x) - kx.\]

On admet que le tableau des variations de la fonction $g$ est le suivant :

\begin{center}
\psset{unit=1cm,arrowsize=2pt 3}
\begin{pspicture}(7.5,2.5)
\psframe(7.5,2.5)\psline(0,1.5)(7.5,1.5)
\psline(1.5,0)(1.5,2.5)
\uput[u](0.75,1.4){$x$} \uput[u](1.65,1.4){$0$} \uput[u](4.5,1.4){$\dfrac 1k$} 
\uput[u](7,1.4){$+ \infty$}
\rput(0.75,0.75){$g(x)$}\uput[u](2,0){$- \infty$}\uput[d](4.5,1.5){$g\left(\frac 1k\right)$}
\uput[u](7,0){$- \infty$}
\psline{->}(2,0.5)(4,1)\psline{->}(5,1)(7,0.5)
\end{pspicture}
\end{center}

	\begin{enumerate}
		\item Donner, en fonction du signe de $g\left(\dfrac 1k\right)$ le nombre de solutions de l'équation $g(x) = 0$.
		\item Calculer $g\left(\dfrac 1k\right)$ en fonction du réel $k$.
		\item Montrer que $g\left(\dfrac 1k\right) > 0$ équivaut à $\ln (k) < - 1$.
		\item Déterminer l'ensemble des valeurs de $k$ pour lesquelles l'équation $\ln (x) = kx$ possède exactement deux solutions.
		\item Donner, selon les valeurs de $k$, le nombre de solutions de l'équation $\ln (x) = kx$.
	\end{enumerate}
\end{enumerate}

\medskip

\end{document}
